\subsection{All-depth local matching at nonzero valuation primes}
\label{sec:ghost-local-rigidity}

\begin{result}{C.1.16}{Proved}{Rational Magnus conjugacy detects pro-\(\ell\) conjugacy}
\label{lem:ghost-magnus-conjugacy}
Let \(F^{(\ell)}=\langle x,y\rangle^{(\ell)}\) be the
free pro-\(\ell\) group, embedded by the Magnus map
\(x\mapsto1+X\), \(y\mapsto1+Y\) in
\(1+(X,Y)\mathbb Z_\ell\langle\!\langle X,Y\rangle\!\rangle\).
The conjugacy relation induced on this group by the full
unit group of
\(\mathbb Q_\ell\langle\!\langle X,Y\rangle\!\rangle\)
is exactly its own pro-\(\ell\) conjugacy relation.
\end{result}
\begin{proof}
For each degree, identify words differing by a cyclic
rotation, and take the sum of the coefficients in each
resulting finite class of words. Together these maps
give the cyclic-word map on formal series. Its values
on \(AB\) and \(BA\) agree: each concatenation of
two words is a cyclic rotation of the concatenation
in the opposite order. All calculations in a fixed
degree are finite. Consequently this map takes the
same value on \(U g U^{-1}\) and \(g\) for every
invertible rational formal series \(U\).

Let \(g,h\in F^{(\ell)}\) have rationally conjugate
Magnus series. Their difference has integral
coefficients, and each of its cyclic coefficient sums
is zero. In each degree it is therefore an integral
linear combination of differences of cyclically
rotated words. Each such difference is a sum of
expressions \(AB-BA\). Thus the truncations of
\(g-h\) have expressions as integral linear
combinations of additive commutators of words.
No rational conjugator is evaluated in an integral
group algebra in what follows.

Let \(P\) be any finite \(\ell\)-group quotient of
\(F^{(\ell)}\). There is a continuous evaluation
homomorphism
\[
 \mathbb Z_\ell\langle\!\langle X,Y\rangle\!\rangle
       \longrightarrow\mathbb Z_\ell[P],\qquad
 X\longmapsto\bar x-1,\quad Y\longmapsto\bar y-1.
\]
Indeed the augmentation ideal of \(\mathbb F_\ell[P]\)
is nilpotent. One proves this by induction on \(|P|\):
choose a central element \(z\) of order \(\ell\),
pass to \(P/\langle z\rangle\), and use
\((z-1)^\ell=0\). If the quotient augmentation ideal
has nilpotence exponent \(N\), the original one has
exponent at most \(\ell N\). Therefore the integral
augmentation ideal \(I(P)\) has a power contained in
\(\ell\mathbb Z_\ell[P]\), and its powers tend to
zero \(\ell\)-adically. This proves convergence of
the evaluation of every integral formal series.

The quotient of \(\mathbb Z_\ell[P]\) by the additive
commutator module is free over \(\mathbb Z_\ell\)
on the conjugacy classes of \(P\). To check the
module explicitly, \(ab-ba\) is the difference of
two conjugate group elements; conversely the
difference of an element and its conjugate is such
an additive commutator. The module is therefore
closed. Evaluation of the integral commutator
truncations above shows that \(\bar g\) and
\(\bar h\) have the same image in this free module.
They must be conjugate in \(P\).

For every open normal subgroup \(N\) of
\(F^{(\ell)}\), let
\[
 C_N=\{v\in F^{(\ell)}:v^{-1}gvh^{-1}\in N\}.
\]
These are nonempty closed subsets of the compact
group \(F^{(\ell)}\). A finite intersection contains
\(C_{\cap N}\), hence is nonempty. Compactness gives
an element in their full intersection, which conjugates
\(g\) to \(h\). The converse follows by applying the
Magnus homomorphism to an actual conjugator. The
injectivity of the Magnus embedding follows from
Lemma~\ref{lem:ghost-hall-model}: a nontrivial element
is detected in some finite \(\ell\)-group quotient,
that quotient has finite nilpotency class, and the
corresponding integral Hall coordinates are detected
by its truncated Magnus series.
\end{proof}

\begin{result}{C.1.17}{Proved}{Singleton local fibre at a boundary valuation}
\label{prop:ghost-local-boundary}
Let \(q\ne\ell\) be primes and let
\(t\in\mathbb Q_q\setminus\{0,1\}\) satisfy
\((v_q(t),v_q(1-t))\ne(0,0)\). In the geometrically
pro-\(\ell\) arithmetic fundamental group of
\(U_{\mathbb Q_q}\), the fibre of the abelianization
map on geometric conjugacy classes of continuous
sections over
\[
             (\kappa_\ell(t),\kappa_\ell(1-t))
\]
consists of exactly the local point-section class of
\(t\). Therefore every all-depth ordinary global
pronilpotent section matches its prescribed ordinary
local point at every coefficient prime \(\ell\ne q\)
and every finite place \(q\) of nonzero valuation
support.
\end{result}
\begin{proof}
We first treat \(m=v_q(t)>0\). The local
\(\ell\)-adic Kummer class of \(1-t\) is zero,
because its principal \(q\)-unit group is uniquely
\(\ell\)-divisible. Work with the genuine action
defined by the tangent \(\overrightarrow{01}\) and
with the generator \(x\) of inertia at zero.

The arithmetic action on the free pro-\(\ell\)
geometric group is unramified at \(q\). Its class-two
action is pure cyclotomic. In the passage from class
\(i\) to class \(i+1\), the kernel of the automorphism
map injects into
\(\operatorname{Hom}(N_1,L_{i+1})\), with Frobenius
weight \(q^i\). Any tame inertia image there is
killed by \(q^i-q\), and is zero for \(i\ge2\);
the module is torsion-free. Wild and non-\(\ell\)
tame inertia also vanish, as their images are in a
pro-\(\ell\) additive group. Induction proves
unramifiedness at every nilpotency class and hence
in the free pro-\(\ell\) group itself. The inverse
limit assertion here follows because every finite
\(\ell\)-group has finite nilpotency class.

Let \(s\) be any section with the specified abelian
class. Its geometric cocycle factors through
\(G_q^{t,\ell}\): on the discarded inertia kernel
the action is trivial, so the cocycle is a homomorphism
from a pro-prime-to-\(\ell\) group to a pro-\(\ell\)
group and is zero. Write \(V\in F^{(\ell)}\) for
its tame inertia image and \(z\in F^{(\ell)}\) for
its geometric Frobenius coordinate. Their relation is
\[
       z\sigma_0(V)z^{-1}=V^q,
\]
where \(\sigma_0\) is the actual tangential Frobenius
automorphism and \(\sigma_0(x)=x^q\). The abelian
coordinates of \(V\) are exactly \((m,0)\);
coboundaries have zero inertia restriction. The
valuation \(m\) is an actual nonzero integer, so
it is nonzero in every \(\mathbb Z_\ell\). A chosen
tame Frobenius and the \(\ell\)-inertia generator
topologically generate \(G_q^{t,\ell}\), whose
projection to \(\widehat{\mathbb Z}\) splits by
the powers of that Frobenius.

Form the rational prounipotent group \(\mathcal U\)
by taking the inverse limit of the rational Hall
groups \(N_i(\mathbb Q_\ell)\). The Hall laws are
rational polynomial identities, so define these
groups over \(\mathbb Q_\ell\); their central
graded terms are the free Lie terms over
\(\mathbb Q_\ell\). The Magnus map embeds them
compatibly in rational formal series. The actual
\(\sigma_0\) extends to \(\mathcal U\) and acts
on its weight-\(i\) term by \(q^i\).

There is a unique \(u\in\mathcal U\) with
\[
                 u^{-1}z\sigma_0(u)=1.
\]
This assertion conjugates the specified Frobenius
section; it does not replace \(\sigma_0\) by a
linearized automorphism. In weight one its solution
is \(u_1=-z_1/(q-1)\). At every later weight \(i\),
changing a lift by a central vector \(w\) changes
the remaining defect by \((q^i-1)w\). This scalar
is nonzero and hence invertible in \(\mathbb Q_\ell\),
giving a unique correction. These corrections are
compatible at successive classes and define \(u\).

Put \(W=u^{-1}Vu\). The tame relation becomes
\(\sigma_0(W)=W^q\). It has abelian coordinates
\((m,0)\). The reference element \(x^m\) has
the same coordinates and satisfies the same equation.
If their logarithms first differed in a weight
\(i\ge2\), the difference in that weight would
be killed by \(q^i-q\), which is nonzero in
\(\mathbb Q_\ell\). Induction gives
\(W=x^m\) in \(\mathcal U\). Thus \(V\) is
rationally formally conjugate to \(x^m\).
Lemma~\ref{lem:ghost-magnus-conjugacy} gives an actual
\(v\in F^{(\ell)}\) with \(v^{-1}Vv=x^m\).

Conjugate the original section by this integral
\(v\). Its Frobenius coordinate \(z'\) satisfies
\(z'x^{qm}(z')^{-1}=x^{qm}\). The centralizer of
\(x^{qm}\) in \(F^{(\ell)}\) is exactly
\(x^{\mathbb Z_\ell}\). Here is an elementary
verification. Over rational formal series, \(X\)
is a formal power series in \((1+X)^{qm}-1\),
since its linear coefficient \(qm\) is nonzero.
An element centralizing \(x^{qm}\) therefore
centralizes \(X\). Comparing coefficients in
\(XA=AX\) shows that every word containing \(Y\)
has zero coefficient: move trailing \(X\)'s to
the beginning until the word ends in \(Y\), where
the equality forces its coefficient to zero.
Thus \(A\in\mathbb Q_\ell[[X]]\). For an
element of \(F^{(\ell)}\), the retraction
\(y\mapsto1\), \(x\mapsto x\) then shows that
it equals \(x^a\) for \(a\in\mathbb Z_\ell\),
using Magnus injectivity. This proves the
centralizer assertion.

It follows that the conjugated section has its
entire geometric cocycle in \(x^{\mathbb Z_\ell}\):
its Frobenius and tame inertia do, and these
topologically generate the tame group. Its class
is completely determined by its first abelian
Kummer class. Indeed the inclusion
\(\mathbb Z_\ell(1)\to F^{(\ell)}\), \(a\mapsto x^a\),
has the first abelian coordinate as an equivariant
retraction, and two such cyclic cocycles with the
same first cohomology class differ by conjugation
by an element of \(x^{\mathbb Z_\ell}\).
Apply the argument also to the actual point section
of \(t\). It has the same first abelian class,
so its section class equals that of \(s\).

For \(v_q(1-t)>0\), the algebraic automorphism
\(t\mapsto1-t\) interchanges the two Kummer
coordinates and reduces to the case just proved.
For the remaining pattern \(v_q(t)=v_q(1-t)<0\),
use \(t\mapsto1/t\), whose coordinate map is
\[
          (A,B)\longmapsto(A^{-1},-BA^{-1}).
\]
Its first valuation is positive and its second
coordinate is the class of \(1-1/t\), a principal
\(q\)-unit. Each automorphism induces an isomorphism
of arithmetic fundamental groups at the transported
base point, and any return base path changes it only
by a geometric inner automorphism. It therefore
induces an actual bijection of the section-conjugacy
fibres under discussion. This proves the statement
for every nonzero ordinary valuation pattern.

Finally a global pronilpotent section projects to
the full geometrically pro-\(\ell\) arithmetic
group and restricts continuously to each local
Galois group. Its abelian class is the given
ordinary pair's local point class. The local
singleton assertion applies componentwise, giving
the final conclusion.
\end{proof}
